\documentclass[a4paper]{IEEEtran}
% \IEEEoverridecommandlockouts
% The preceding line is only needed to identify funding in the first footnote.
% If that is unneeded, please comment it out.
\usepackage[
    backend=biber,
    style=ieee
]{biblatex}
\DeclareFieldFormat[article]{volume}{\mkbibbold{#1}}
\usepackage{amsmath,amssymb,amsfonts}
\usepackage{graphicx}
\usepackage{textcomp}
\usepackage{xcolor}
\def\BibTeX{{\rm B\kern-.05em{\sc i\kern-.025em b}\kern-.08em
    T\kern-.1667em\lower.7ex\hbox{E}\kern-.125emX}}

\usepackage[hidelinks]{hyperref}
\usepackage{cleveref}
\usepackage{booktabs}
\usepackage{algorithm}
\usepackage{algpseudocode}

\usepackage{listings,newtxtt}
\lstset{
    basicstyle=\footnotesize\ttfamily,
    keywordstyle=\bfseries,
    frame=leftline,
    framesep=4pt,
}
\crefname{lstlisting}{listing}{listings}
\Crefname{lstlisting}{Listing}{Listings}
\usepackage{bm}
\usepackage{tabularx}

\addbibresource{references.bib}

\begin{document}

\title{TODO: Project 2 Title}

\author{
    \IEEEauthorblockN{Sigurd Riseth}\\
    \IEEEauthorblockA{\textit{Department of Physics}\\
    \textit{University of Oslo}\\
    Oslo, Norway \\
    sigurri@uio.no}
}

\maketitle

% Abstract. 120-180 words, one paragraph, no citations or equations.
% Past tense for what you did, and at least one concrete result.
\begin{abstract}
TODO: abstract.
\end{abstract}

\begin{IEEEkeywords}
TODO, keywords
\end{IEEEkeywords}

%\tableofcontents

\section{Introduction}\label{sec:introduction}

% An introduction where you explain the aims and rationale for the physics case
% and what you have done. At the end of the introduction you should give a brief
% summary of the structure of the report.
%
% State explicitly what is your own code and what is scikit-learn/other
% libraries. End with a paragraph on the report's structure and the repo link.
% No results in the introduction.

TODO: introduction.

\section{Methods}
\label{sec:method}

% Theory, methods, algorithms. Assume the reader has familiarity with basic
% linear algebra and statistics. Break into \subsection{}s as the content
% develops, e.g. one per model/algorithm.

TODO: methods.

\section{Implementation}
\label{sec:implementation}

% Describe the data, the preprocessing, the structure of the code, and how it
% was verified (unit tests, comparison against scikit-learn/known solutions,
% etc.). Link to the repository.

TODO: implementation.

\cite{riseth_2026a}

\subsection{Data and preprocessing}
\label{subsec:data}

TODO: data and preprocessing.

\section{Results and Discussion}
\label{sec:results-discussion}

% Present and discuss results together, ordered to match the structure of
% \cref{sec:method}. Break into \subsection{}s as the content develops.

TODO: results and discussion.

% \begin{figure}[htbp]
%     \centering
%     \includegraphics[width=\linewidth]{figures/PLACEHOLDER.pdf}
%     \caption{TODO.}
%     \label{fig:placeholder}
% \end{figure}

\section{Conclusion}\label{sec:conclusion}

% Summarize the central findings, state the limitations of the study, and
% outline concrete future work/extensions.

TODO: conclusion.


\printbibliography

\appendix

% Additional calculations used to validate the codes
% Selected calculations, these can be listed with few comments
% Listing of the code if you feel this is necessary

% You can consider moving parts of the material from the methods section
% to the appendix or your GitHub/Gitlab repository. You can also place
% additional material on your webpage.

\subsection{Use of AI/LLM tools}
\label{app:ai-tools}

Large language models were used throughout this project, following the course
guidelines\footnote{\url{https://github.com/EducationalMaterialUiO/MachineLearningUiO/blob/main/LLM_Usage_Declaration_Guidelines.md}},
which also define the contribution levels below. All code was tested, all
derivations were checked, and every claim in the report was read, understood
and verified by the author.

\subsubsection{Tools used}
\begin{itemize}
    \item TODO: tools, models, and dates.
\end{itemize}

TODO: summary of how each tool was used.

\subsubsection{Text writing and editing}
\Cref{tab:llm-writing-contribution} gives the contribution level for each part
of the report.

\begin{table}[htbp]
    \centering
    \caption{LLM contribution to text writing and editing (levels 0--3).}
    \label{tab:llm-writing-contribution}
    {\footnotesize
    \begin{tabularx}{\columnwidth}{@{}lcX@{}}
        \toprule
        \textbf{Section} & \textbf{Level} & \textbf{Notes} \\
        \midrule
        Abstract & TODO & \\
        Introduction & TODO & \\
        Methods & TODO & \\
        Implementation & TODO & \\
        Results and discussion & TODO & \\
        Conclusion & TODO & \\
        Appendices & TODO & \\
        \bottomrule
    \end{tabularx}}
\end{table}

\subsubsection{Code generation and assistance}
\Cref{tab:code-generation} gives the level for each source file (levels 0--4).
Functions with level 2 or above carry an \texttt{LLM-assisted} tag in their
docstring.

\begin{table*}[htbp]
    \centering
    \caption{LLM contribution to the code (levels 0--4).}
    \label{tab:code-generation}
    {\footnotesize
    \begin{tabularx}{\textwidth}{@{}l c X@{}}
        \toprule
        \textbf{File} & \textbf{Level} & \textbf{Description} \\
        \midrule
        TODO & TODO & \\
        \bottomrule
    \end{tabularx}}
\end{table*}

\subsubsection{Declared and undeclared use}
TODO: one example of a use that must be declared, and one example of a use
that need not be declared (e.g. using an LLM as a tutor and then writing
about it in your own words).

\subsection{Additional results}
\label{app:additional-results}

% Supplementary figures/tables that support but don't belong in the main
% results section.

TODO: additional results.


\end{document}
